\documentclass[11pt]{article}
\usepackage{amsmath,amssymb,amsthm,mathtools}
\usepackage{bm}
\usepackage{geometry}
\usepackage{booktabs}
\usepackage{array}
\usepackage{tabularx}
\usepackage{graphicx}
\usepackage{hyperref}
\usepackage{enumerate}

\geometry{letterpaper, margin=1in}

\newcommand{\boldI}{\mathbf{I}}
\newcommand{\boldG}{\mathbf{G}}
\newcommand{\boldC}{\mathbf{C}}
\newcommand{\vr}{\mathbf{r}}
\newcommand{\vE}{\mathbf{E}}
\newcommand{\vmu}{\boldsymbol{\mu}}
\newcommand{\valpha}{\boldsymbol{\alpha}}
\newcommand{\boldT}{\mathbf{T}}
\DeclareMathOperator*{\argmax}{arg\,max}

\title{\textbf{Formation Engine Methodology} \\
\Large Final Statement --- Master's Thesis Conclusion}
\author{Formation Engine Development Team}
\date{Version 3.1.0}

\begin{document}
\maketitle

\tableofcontents
\newpage

% ═══════════════════════════════════════════════════════════════════════════
\section{Restatement of the Thesis}
% ═══════════════════════════════════════════════════════════════════════════

This work set out to answer a single, concrete question:

\begin{quote}
\textit{Can a deterministic, classical atomistic engine generate
physically admissible molecular and crystalline structures from
elemental identity alone, without reference to molecular databases,
pre-solved crystal structures, or heuristic geometry templates ---
and can it do so in a manner that is fully reproducible,
explicitly auditable, and scientifically credible?}
\end{quote}

The answer established over thirteen methodology sections, eleven
validation suites comprising 1,013 tests, and an autonomous
self-audit infrastructure is: \textbf{yes, within declared boundaries}.

The Formation Engine is that instrument.  It is not a prediction
claim.  It is not a materials-discovery oracle.  It is a reproducible
computational tool for generating and evaluating candidate atomistic
structures under explicit physical assumptions, targeted at the
intermediate regime of 100--10,000 atoms where quantum methods
become computationally intractable and continuum models must
assume rather than derive structure.

% ═══════════════════════════════════════════════════════════════════════════
\section{Summary of Contributions}
% ═══════════════════════════════════════════════════════════════════════════

The contributions of this work are organized by layer.

\subsection{Scientific Problem Identification (\S1)}

The foundational contribution is a precise diagnosis of the
\textbf{emergence boundary} problem.  The gap between quantum
chemistry (limited to $\sim$200 atoms at DFT level) and bulk
continuum models is not merely a computational one.  It is a
\textit{conceptual} one: existing frameworks have no natural
language for structure as a computational output rather than
as a starting assumption.

The formation engine resolves this by treating initialization
not as a preparatory step but as a \textbf{primary experimental
parameter}.  Every structure generated by the engine carries
full provenance: the random seed, the temperature and density
boundary conditions, the force field parameters, and the
convergence criteria.  This converts initialization from an
opaque procedural artifact into an auditable scientific record.

\subsection{Canonical State Representation (\S0, \S2)}

A complete, partition-respecting state representation was
developed:
\[
\mathcal{S} = (N,\,\mathbf{M},\,\boldsymbol{\tau},\,\mathbf{X},\,
\mathbf{V},\,\mathbf{Q},\,\mathbf{B},\,\mathbf{F},\,\mathcal{E},\,
\mathcal{L},\,\mathcal{B})
\]

partitioned into three physically motivated tiers:
\begin{enumerate}
\item \textbf{Identity} $(N, \mathbf{M}, \boldsymbol{\tau})$ ---
      immutable across the simulation lifetime.
\item \textbf{Phase} $(\mathbf{X}, \mathbf{V}, \mathbf{Q}, \mathbf{B})$
      --- governed by equations of motion; updated by the integrator.
\item \textbf{Scratch} $(\mathbf{F}, \mathcal{E}, \mathcal{L}, \mathcal{B})$
      --- fully determined by Phase; recomputed, never saved as
      independent state.
\end{enumerate}

This partition is not cosmetic.  It enforces a discipline that prevents
common simulation errors: reading forces before they are evaluated,
treating energy as an independent degree of freedom, or accumulating
numerical drift in quantities that should be recomputed from scratch.

The per-particle identity vector
$\boldI_i = (Z_i, A_i, Q_i, \Sigma_i, \Lambda_i, \Theta_i)$
grounds every interaction parameter in a single periodic table
lookup, satisfying the \textbf{periodic table as sole authority}
axiom.

\subsection{Canonical Particle Container and Quantum Descriptors}

The identity--state decomposition was extended in the unified
methodology to a full canonical container representation:
\[
\mathcal{C}_i = \bigl(\boldI_i,\;\boldG_i,\;\boldC_i^{\,\mathrm{comp}}\bigr)
\]
where $\boldG_i = (s_i, c_i, f_i, \kappa_i, \eta_i)$ is the
quantum descriptor bundle (spin, colour charge, flavour,
projection coupling, information density) and
$\boldC_i^{\,\mathrm{comp}} = (\Gamma_i, \Xi_i, \Pi_i, \Upsilon_i)$
is the computational annotation payload carrying complexity class,
kernel selector, solver priority, and annotation confidence.

The Container Invariance Principle states:
$\mathcal{C}_i \equiv \mathcal{C}_j$ iff
$\boldI_i = \boldI_j \land \boldG_i = \boldG_j$;
computational layers may differ.
This formalism ensures that particle identity is
physics-determined, while computational treatment remains
adaptive.

The framework was further extended to cell and world
containers.  A simulation cell is
$\mathcal{K} = (\mathbf{K}^{\mathrm{id}},\,
\mathbf{K}^{\mathrm{geom}},\,\mathbf{K}^{\mathrm{comp}})$
with geometry defined by a $3 \times 3$ lattice matrix
$\mathbf{L} = [\mathbf{a}\;\mathbf{b}\;\mathbf{c}]^\top$,
$V = \det(\mathbf{L})$.
A world (multi-cell simulation instance) is
$\mathcal{W} = (W^{\mathrm{id}},\, W^{\mathrm{top}},\,
W^{\mathrm{ctrl}},\, W^{\mathrm{comp}})$
with a connectivity graph $G_{\mathcal{W}} = (V, E)$ between cells.

\subsection{Interaction Model (\S3)}

The nonbonded potential is a Lennard-Jones 12-6 interaction:
\[
U_{\mathrm{LJ}}(r_{ij}) = 4\varepsilon_{ij}
\left[\left(\frac{\sigma_{ij}}{r_{ij}}\right)^{12}
- \left(\frac{\sigma_{ij}}{r_{ij}}\right)^{6}\right]
\cdot S(r_{ij};\,r_{\mathrm{on}},\,r_c)
\]
with Lorentz-Berthelot combining rules
($\varepsilon_{ij} = \sqrt{\varepsilon_i\varepsilon_j}$,
$\sigma_{ij} = (\sigma_i + \sigma_j)/2$) and a quintic switching
function active over $[0.9\,r_c,\,r_c]$ at $r_c = 10$\,\AA.

Bonded terms (harmonic bond, angle, torsion) are topology-aware
and inferred directly from atomic positions via covalent radius
matching with tolerance $f = 1.2$.  This makes the bond graph a
\textit{derived} quantity of the geometry, not an independent
input, consistent with the no-hidden-state requirement.

All parameters ($\varepsilon$, $\sigma$, $r_{\mathrm{cov}}$, $r_{\mathrm{vdW}}$)
are drawn exclusively from UFF (Universal Force Field), covering
$Z = 1$--$102$.

\subsection{Thermodynamics and Integration (\S4, \S5)}

The canonical ensemble is maintained via a velocity-rescaling
thermostat, implemented as a \textit{declarative} physical model
rather than a numerical stabilizer.  The thermostat's effect on
the energy budget is tracked explicitly in the energy ledger
$\mathcal{E}$.  The Nos\'e-Hoover chain extension (\S4) provides
a path to rigorously canonical sampling for future work.

Time integration uses the \textbf{Velocity Verlet} algorithm with
default timestep $\Delta t = 1$\,fs.  Numerical stability was
verified through total-energy conservation tests with Coulomb
disabled.  The integrator produces bit-identical trajectories
given identical initial state and seed, satisfying the
\textbf{deterministic core} axiom.

Structure minimization uses the \textbf{FIRE} (Fast Inertial
Relaxation Engine) algorithm, which was chosen for its compatibility
with the existing Velocity Verlet integrator and its empirical
convergence reliability on systems with 2--10,000 atoms.
Convergence criteria are:
$\max_i |\mathbf{F}_i| < 10^{-6}$\,kcal/(mol\,\AA) and
$|\Delta E|/N < 10^{-10}$\,kcal/mol.

\subsection{Formation Physics and Statistical Classification (\S6, \S7)}

The formation pipeline transforms a seed and a molecular formula
into a classified structure through a reproducible five-stage
protocol: initialization, high-temperature exploration, staged
annealing, FIRE minimization, and structural classification.

The classification system assigns every formation outcome to one
of six categories: \texttt{stable}, \texttt{metastable},
\texttt{unstable}, \texttt{fragment}, \texttt{collapsed},
or \texttt{timeout}.  Classification uses quantitative metrics ---
energy per atom, maximum residual force, bond-length deviation
from reference, and coordination-number consistency --- with
explicit numerical thresholds, not heuristic labels.

Statistical analysis across seed ensembles (\S7) provides
formation probability, stability score, and structural diversity
index.  These are the quantities that make the engine useful as
a \textit{screening} tool: running 100 seeds for a given formula
produces a statistically meaningful picture of its accessible
configuration space.

\subsection{Reaction and Electronic Extensions (\S8, \S8b, \S9)}

Section~8 documents the reaction engine: a graph-rewrite system
that proposes, filters, and applies bond-topology changes under
thermodynamic gating.  Section~8b introduces heat-gated reaction
control, in which reaction proposals are accepted or rejected
based on local thermal energy rather than global temperature,
enabling spatially heterogeneous reaction landscapes.

Section~9 documents the electronic state extension: partial
charges via electronegativity equalization, polarizability
corrections, and the interface to the (currently disabled)
Coulomb term.  The Coulomb implementation is structurally
complete; re-enabling it requires only a dedicated stability
audit of the integrator-Coulomb coupling.

\subsection{Polarization Models}

The unified methodology introduces self-consistent field (SCF)
induced dipole polarization to address the 10--30\% underestimation
of hydrogen-bond energies and 50\%+ underestimation for metal
complexes inherent in fixed-charge models.  Each particle acquires
an induced dipole:
\[
\vmu_i = \valpha_i \cdot \vE_i^{\text{loc}}
\]
where the local field includes both permanent and induced
contributions coupled through the dipole--dipole interaction tensor
$\boldT_{ij}$.  The resulting $3N \times 3N$ linear system
$(\mathbf{I} - \mathbf{A})\boldsymbol{\mu} = \vE^{\text{perm}}$
is solved iteratively.  Thole damping with parameter $a \approx 2.6$
prevents the polarization catastrophe at short range.  Atomic
polarizabilities are drawn from Miller (1990) tabulations,
maintaining the periodic-table-as-sole-authority axiom.

\subsection{Environment-Responsive Beads}

The coarse-grained layer was extended with environment-responsive
beads that transform from static descriptor carriers into active
material participants.  Each bead carries an extended state:
\[
\mathbb{B}_B = \bigl\{
  \underbrace{\mathbf{R}_B,\,\mathbf{Q}_B,\,\{c_{\ell m}^{(k)}\},\,M_B,\,Q_B^{\mathrm{tot}},\,\boldI_B}_{\text{intrinsic}},\;
  \underbrace{\rho_B,\,C_B,\,P_{2,B},\,\eta_B}_{\text{environment}}
\bigr\}
\]

Three fast observables (local density $\rho_B$, coordination number
$C_B$, orientational order parameter $P_{2,B}$) are recomputed every
dynamics step.  A slow internal memory variable $\eta_B \in [0,1]$
evolves via first-order relaxation:
\[
\frac{d\eta_B}{dt} = \frac{1}{\tau}\left(
  \alpha\,\hat{\rho}_B + \beta\,\hat{P}_{2,B} - \eta_B
\right)
\]
where $\tau$ is the memory timescale and $\alpha + \beta = 1$.
Environment state couples back into interaction kernels via
$K_k(\ell, r;\,\eta_A, \eta_B) = K_k(\ell, r) \cdot
(1 + \gamma_k\,\bar{\eta}_{AB})$, enabling emergent phenomena
including spontaneous stacking, electrostatic screening, hysteresis,
and nucleation asymmetry.

\subsection{Library Reactions and Thermodynamic Verification}

Three reference reactions were implemented with full Hess's-law
thermodynamics and mass-balance verification:
\begin{enumerate}
\item \textbf{Benzene nitration} (electrophilic aromatic substitution):
      $\Delta H_{\mathrm{rxn}} = -35.56$\,kcal/mol.
\item \textbf{Thorium oxalate precipitation} (ligand exchange,
      product phase: sludge).
\item \textbf{Copper nitrate decomposition} (thermal,
      $\Delta H_{\mathrm{rxn}} = +103.72$\,kcal/mol, endothermic).
\end{enumerate}
Each reaction carries verified atom conservation through a
mass-balance system that checks
$\forall E: \sum_i \nu_i n_{E,i} = \sum_j \nu_j n_{E,j}$
with tolerance $10^{-6}$.

\subsection{BeadFIRE Visualization and Reporting Pipeline}

A complete automated outcome-reporting and visualization pipeline
was delivered for BeadFIRE (Lattice-Level FIRE) minimization.
The pipeline comprises:
\begin{itemize}
\item C++ position recording at every FIRE step (stored in
      \texttt{BeadFIREStep::positions})
\item Structured Markdown report generation with convergence
      metrics, energy decomposition, and \LaTeX\ equations
\item Machine-readable CSV trajectory export with full bead
      position history
\item Python visualization producing publication-quality
      (300\,DPI) XY and XZ plane trajectory projections, energy
      convergence plots, force decay curves, and FIRE parameter
      evolution
\end{itemize}
All data flows from the actual FIRE kernel with no synthetic
overlays, maintaining anti-black-box integrity.  The system is
32-bit compatible and portable across platforms.

\subsection{Multiscale Projection and Future Architecture (\S10, \S12, \S13)}

Section~10 establishes the four-tier scale hierarchy
(molecular, cluster, bulk, coarse-grained) and defines the
supercell construction protocol.  Supercell geometry is replicated
from the unit cell using:
\[
\mathbf{r}_{i,pqr} = \mathbf{r}_i + p\,\mathbf{a}
+ q\,\mathbf{b} + r\,\mathbf{c}
\]
with bond graphs re-inferred post-construction from atomic
geometry.  This ensures that topology is always consistent with
the actual geometry, eliminating a common source of silent
errors in periodic structure workflows.

Sections~12 and~13 document the extension architecture: the
interfaces for machine-learning potential correction, reactive
bond-order dynamics, triclinic cell support, and variable-cell
relaxation.  Each extension is specified at the interface level
with explicit fallback requirements.  The classical core remains
independently evaluable regardless of which extensions are active.

\subsection{Deterministic Pipeline and Self-Audit (\S11)}

The self-audit infrastructure is the enforcement mechanism for
the reproducibility guarantee.  It comprises three tools:

\begin{enumerate}
\item \textbf{Failure classifier} --- autonomously categorizes
      outcomes from large run batches ($>$10,000 formations) into
      the six formation classes, producing structured JSON records.
\item \textbf{Gap targeter} --- analyzes the coverage of parameter
      space and generates targeted new runs to fill identified gaps,
      enabling systematic exploration rather than random sampling.
\item \textbf{Regression detector} --- compares classification
      outcomes across parameter versions, flags any molecule
      whose formation class or stability score changed, and
      enforces a binary verdict: \textsc{no regressions} or
      \textsc{review required}.
\end{enumerate}

The reproducibility contract is strict: given
(formula, seed, parameters), the engine produces \textbf{bit-identical
output} across platforms, compiler versions, and run sequences.
This is enforced by seed-controlled \texttt{std::mt19937} sampling,
index-ordered force accumulation, and fixed-point summation for
energy aggregation.

% ═══════════════════════════════════════════════════════════════════════════
\section{Empirical Results}
% ═══════════════════════════════════════════════════════════════════════════

\subsection{Formation Determinism}

The unified methodology establishes the Formation Determinism
Postulate as a formal result:
\[
\boxed{\;M(t^{*}) = \mathcal{F}\!\bigl(\mathcal{E}\bigr)\;}
\]
where the formation condition vector
$\mathcal{E} = [T(t),\,c(t),\,p(t),\,P,\,R]$ captures thermal
history, composition history, pressure history, pathway class, and
environmental regime.  Structure is a deterministic function of
formation conditions --- not an independent variable, not an
assumption, but a computed output with full provenance.

\subsection{Variable Hierarchy Reduction}

Across the 1,013-test validation suite, the seven reported state
variables $(\rho, C, P_2, \hat{\rho}, \hat{P}_2, f, \eta)$ were
shown to reduce to four independent degrees of freedom:
\[
\text{Independent:}\quad \rho,\; C,\; P_2,\; \eta
\qquad
\text{Derived:}\quad \hat{\rho},\; \hat{P}_2,\; f
\]
This reduction simplifies the state space of the environment-responsive
layer by nearly half, making parameter sweeps tractable and
confirming that the derived quantities carry no independent
information.

\subsection{Causal Chain Confirmation}

The empirically confirmed causal structure is:
\[
\underbrace{\text{geometry}}_{\text{scene}}
\;\longrightarrow\;
\underbrace{\rho,\; C,\; P_2}_{\text{fast obs.}}
\;\longrightarrow\;
\underbrace{f}_{\text{target}}
\;\longrightarrow\;
\underbrace{\eta}_{\text{memory}}
\]
Geometry determines fast observables.  Fast observables determine the
driving force $f$.  The driving force determines the memory state
$\eta$.  No reverse causation was observed.  The system exhibits
pure first-order relaxation with no multistability.

\subsection{Validation Scale}

The validation infrastructure grew from the initial 35 Phase-1
audit checks to 1,013 tests across eleven suites:

\begin{center}
\begin{tabular}{lrl}
\toprule
Suite & Tests & Focus \\
\midrule
Suite \#1  & 60  & Channel kernels, SH rotation \\
Suite \#2  & 77  & Observable correctness, $\eta$ dynamics \\
Track 2    & 39  & Monte Carlo robustness (4000+ trials) \\
Suite \#3  & 52  & Large-$N$ formation studies \\
Suite \#4  & 113 & Dynamic regimes, invariance \\
Suite \#5  & 131 & Ensemble proxies, response validation \\
Suite \#6  & 85  & Macro property precursors \\
Suite \#7  & 109 & Property learning pipeline \\
Suite \#8  & 193 & Calibration, OOD detection \\
CLI Suite  & 72  & Integration tests \\
Bead Layers & 82 & Visualization data model \\
\midrule
\textbf{Total} & \textbf{1,013} & All passing \\
\bottomrule
\end{tabular}
\end{center}

Key empirical findings across the suite:
\begin{enumerate}
\item Large-$N$ runs ($N = 64, 125, 216$) are numerically stable.
\item Edge vs.\ bulk differentiation emerges naturally from the
      environment-responsive layer.
\item The causal chain geometry $\to$ observables $\to$ target $\to$
      $\eta$ is empirically confirmed without exception.
\item The system exhibits pure first-order relaxation with no
      multistability.
\item Monte Carlo robustness testing over 4,000+ trials confirms
      statistical stability.
\end{enumerate}

% ═══════════════════════════════════════════════════════════════════════════
\section{What the Four Axioms Enabled}
% ═══════════════════════════════════════════════════════════════════════════

The formation engine was built on four design axioms (\S1).  In
retrospect, each axiom had a specific payoff.

\paragraph{Explicit units everywhere.}
Eliminating reduced units forced every quantity to carry a physical
meaning.  This caught several parameter inconsistencies during
development that would have been silent errors in a reduced-unit
framework.  It also makes outputs directly interpretable without
a unit-conversion appendix.

\paragraph{No silent domain switching.}
Declaring the force field, integrator, and boundary conditions at
session initialization made the simulation state globally inspectable
at any point.  There is no ``which model is currently active?''
ambiguity.  This reduced debugging time significantly and made
the regression detector straightforward to implement.

\paragraph{Deterministic core with stochastic exploration.}
Separating physics (deterministic) from exploration (stochastic,
seed-controlled) made ensemble analysis scientifically meaningful.
If two seeds produce different structures, that difference is
\textit{physical} (distinct metastable states), not numerical
artifact.  This distinction is essential for using the engine
as a screening tool.

\paragraph{Periodic table as sole authority.}
The constraint to derive all parameters from $Z$ alone eliminated
an entire class of data-dependency problems.  The engine can in
principle evaluate any combination of elements up to $Z = 102$
without additional parameterization effort.  This generality
comes at the cost of accuracy for specific systems (UFF is not
a high-fidelity force field), but that cost is \textit{declared}
and \textit{bounded}.

% ═══════════════════════════════════════════════════════════════════════════
\section{Declared Limitations}
% ═══════════════════════════════════════════════════════════════════════════

Scientific honesty requires explicit enumeration of what this work
does not accomplish.

\subsection{Coulomb Interactions Are Disabled}

The Coulomb term is implemented but disabled at the kernel level
(\texttt{U\_coul = 0.0}, \texttt{F\_coul\_r = 0.0} in
\texttt{lj\_coulomb.cpp}).  This was triggered by an observed
integrator-Coulomb coupling instability in NaCl dynamics.  The
root cause is identified as a timestep-dependent energy drift
in the force accumulation path, but has not been resolved.

Consequence: all polar and ionic systems are currently evaluated
using van der Waals interactions only.  Relative energies for
systems where electrostatics dominate (salts, charged biomolecules,
strongly polar solvents) are not reliable.  This does not affect
noble gases, nonpolar organics, or metallic systems evaluated
at the LJ level.

\subsection{Orthogonal Boxes Only}

Periodic boundary conditions are implemented for orthogonal boxes
only.  Non-cubic crystals (monoclinic, triclinic) require the
full $3 \times 3$ lattice tensor $\mathbf{h}$ and its inverse.
This extension is architecturally specified in \S10 but not
implemented.

\subsection{Static Bond Graph}

The bond graph $\mathbf{B}$ is inferred at initialization and
does not update during dynamics.  This is appropriate for rigid
molecules and non-reactive materials, but it means bond-breaking
and bond-formation events are not captured during production MD.
The reactive dynamics extension (discrete graph-mutation rules
and ReaxFF-style bond-order potentials) is architecturally
specified in \S13 but not implemented.

\subsection{Coarse-Grained Layer Not Implemented}

The xyzC format reserves slots for coarse-grained data
(bead mapping, PMF reference, bulk density), but the
atomistic-to-coarse-grained mapping is not implemented.
All reserved slots are explicitly null in current output.
They must never be interpreted as computed values.

\subsection{Force Field Accuracy}

UFF provides broad elemental coverage with moderate accuracy.
It was not designed for high-precision energetics and is not
competitive with system-specific force fields (OPLS, AMBER, CHARMM)
for biomolecular or organic-chemical applications.  The engine
should be understood as a \textit{structure generation} and
\textit{screening} instrument, not a quantitative thermodynamic
predictor.  Formations identified as stable candidates by the
engine are appropriate inputs for refinement with higher-fidelity
methods (DFT, CCSD(T)), not for direct comparison to experimental
energies.

% ═══════════════════════════════════════════════════════════════════════════
\section{Forward Path}
% ═══════════════════════════════════════════════════════════════════════════

The following extensions are prioritized by their scientific impact
relative to implementation effort.

\subsection{Immediate: Coulomb Stability Audit}

Re-enabling the Coulomb term is the single highest-impact near-term
task.  The required steps are:
\begin{enumerate}
\item Isolate the integrator-Coulomb coupling instability
      in a minimal NaCl dimer testcase.
\item Apply a standard fix (reduced timestep, Wolf summation
      truncation, or Ewald splitting) and confirm energy
      conservation in NVE dynamics.
\item Pass the Phase~2 Coulomb stability audit.
\item Uncomment the two disabled lines in \texttt{lj\_coulomb.cpp}.
\end{enumerate}
This unlocks the full intended domain of the engine.

\subsection{Near-Term: Triclinic Cell Support}

Replacing the orthogonal box representation with the full
lattice tensor $\mathbf{h}$ enables simulation of all
14 Bravais lattices.  This is required for the engine to
function as a genuine crystalline structure generator.
The extension is architecturally specified; implementation
effort is estimated at one to two weeks.

\subsection{Near-Term: Visualization Frontend}

A dedicated atomistic visualization command
(\texttt{viz ``molecule name'' --radius}) is the current
active development direction.  The BeadFIRE visualization
pipeline (Markdown + CSV + Python figures) demonstrates
the reporting pattern; extending it to the atomistic layer
connects the engine's outputs directly to an inspectable,
figure-ready interface, closing the loop between computation
and research reporting.

\subsection{Near-Term: SCF Polarization Integration}

The polarization model (SCF induced dipoles with Thole damping)
is fully specified and parameterized.  Integrating it into the
production force loop re-enables accurate treatment of
hydrogen-bonded and metal-coordination systems.  The $3N \times 3N$
linear solve adds computational cost but is bounded by the
existing cutoff infrastructure.

\subsection{Mid-Term: Dynamic Bond Graph}

Implementing discrete graph-mutation rules for bond
formation and breaking converts the engine from a
structure-relaxation tool into a genuine reaction-exploration
tool.  This is the enabling step for the heat-gated
reaction-control system documented in \S8b.

\subsection{Mid-Term: Atomistic to Coarse-Grained Mapping}

Implementing the CG mapping protocol populates the currently
null slots in the xyzC format and enables the engine to serve
as a bottom-up parameterization source for coarse-grained models,
completing the intended cross-scale pipeline.

\subsection{Long-Term: ML Potential Correction}

The \texttt{MLCorrection} interface specified in \S13 provides a
structurally clean extension point for neural network potential
corrections.  The additive correction architecture
$U_{\text{total}} = U_{\text{UFF}} + \Delta U_{\text{ML}}$
ensures that the classical core remains independently evaluable
and that ML correction is strictly optional.

% ═══════════════════════════════════════════════════════════════════════════
\section{Situating the Work}
% ═══════════════════════════════════════════════════════════════════════════

The Formation Engine occupies a specific and intentionally
narrow position in the computational materials science landscape.

It is \textbf{not} a replacement for DFT.  Quantum accuracy requires
quantum methods.  The engine is intended to generate physically
plausible candidates for DFT refinement, not to substitute for it.

It is \textbf{not} a machine-learning surrogate.  It does not learn
from data and does not generalize from a training set.  It applies
declared physics to declared inputs, making its domain and its
failure modes explicit.

It is \textbf{not} a molecular dynamics package in the traditional
sense.  LAMMPS, GROMACS, and NAMD are equilibration engines
optimized for production sampling.  The Formation Engine is a
\textit{structure discovery} engine optimized for reproducible
candidate generation.

What the engine \textit{is}: a deterministic, auditable atomistic
instrument for generating and classifying candidate structures in
the intermediate size regime, grounded exclusively in the periodic
table, and designed to serve as the first stage of a multi-fidelity
computational investigation pipeline.  With the unified methodology
extensions --- canonical particle containers, environment-responsive
beads, SCF polarization, and automated reporting --- the engine now
spans from quantum descriptor bundles at the particle level to
research-quality figure generation at the output level.

The anti-black-box design philosophy --- every mapping decision,
every parameter, every intermediate result inspectable and traceable
--- is not a stylistic preference.  It is the scientific requirement
that makes the engine's outputs usable in a research context where
results must be understood, not merely reproduced.  The BeadFIRE
visualization pipeline, which produces publication-quality XY and XZ
plane projections from actual kernel trajectory data with no
synthetic overlays, is a concrete expression of this philosophy:
what the engine computes is what the researcher sees.

% ═══════════════════════════════════════════════════════════════════════════
\section{Final Statement}
% ═══════════════════════════════════════════════════════════════════════════

The work presented in this thesis establishes that a deterministic
classical atomistic engine can generate, classify, and audit molecular
and crystalline candidate structures starting from elemental identity
alone, with no reference to experimental databases, heuristic templates,
or opaque initialization procedures.

Thirteen methodology sections, eleven validation suites comprising
1,013 tests, three reference reactions with verified thermodynamics,
an environment-responsive bead layer with empirically confirmed
causal structure, a complete automated reporting and visualization
pipeline, and an autonomous self-audit infrastructure constitute a
functioning scientific instrument with fully declared scope, fully
declared limitations, and a clear forward path.

The formation engine does not claim to solve materials discovery.
It claims something more modest and more durable: \textbf{it makes
the intermediate-scale structure generation step explicit, reproducible,
and scientifically legible}.

The Formation Determinism Postulate ---
$M(t^{*}) = \mathcal{F}(\mathcal{E})$ ---
is the formal distillation of this claim.  Structure is not assumed;
it is computed.  Formation conditions are not hidden; they are
recorded.  Results are not approximate; they are bit-identical.

In a field where initialization is routinely treated as a procedural
detail, treating it as a primary experimental parameter is not a
minor engineering choice.  It is a statement about what computational
science owes to reproducibility.

The engine exists to make that statement in working code.

\vspace{2em}
\noindent\rule{0.4\textwidth}{0.4pt}

\noindent\textit{Formation Engine v3.1.0 ---
13 sections, 1,013 validation tests across 11 suites,
Phase~1 audit passed, all tests passing.}

\end{document}
